\section{Discussion}
\label{sec:discussion}

\subsection{Implications}
If the one-versus-one ladder generalises, then isolated-agent
evaluations will systematically understate operational risk. A
deployment that ``passed'' a static red team can still be unsafe once
the adversary is allowed to remember last week's detections. Conversely,
a defender that looks weak against a noisy scripted attacker may be
adequate against a static one and inadequate against a persistent one.
Certification of agentic systems should therefore state the adaptation
level of both sides, not only the tools on the system under test.

The emergence results caution against a second mistake: treating every
unexpected tool sequence as a new vulnerability class. Delayed attack
is not a new capability; it is a new \emph{use} of $\mathtt{wait}$.
Governance and monitoring should watch strategy signatures (CEP, novelty)
as well as individual actions.

\subsection{Security as an Evolutionary Process}
Equation~\eqref{eq:security-fn} is the intended conceptual contribution.
The practical consequence is that defence cannot be a one-time policy
search. It has to be a standing co-evolutionary process: defenders
need persistent memory at least as long as the attackers they face, and
they need metrics that detect when the opponent has moved
(CEP)~\cite{tambe2011security,hofbauer2003evolutionary}.
Deception versus detection versus counter-deception is the loop we
expect to dominate once LLM agents replace the heuristic policies of
Section~\ref{sec:results}.

\subsection{Limitations}
The present evidence is a heuristic ladder on a simulated network.
Language-model policies may discover different strategies, including
ones that exploit the prompt rather than the environment; they may also
fail to exploit the environment at all. The simulator's detection
physics is a design choice ($\lambda$-exponential in suspicion); a
different IDS model could change the value of stealth. Partial
observability is real but coarse. The Kubernetes backend has been
validated as an interface-compatible drop-in, not yet as a
provisioned range with live NetworkPolicies. Population experiments
(E4--E8) may reverse qualitative conclusions: more defenders can help
or, past a critical interaction density, create new attack paths
through information propagation. Finally, CEP uses a hand-crafted
signature; a learned embedding of trajectories would be a natural
refinement.

\subsection{Threats to Validity}
Internal: a single seed for the ladder; we will report confidence
intervals over seeds in the camera-ready. Construct: ASR/DR do not
capture partial compromise. External: a three-host toy network is not
an enterprise fleet. We argue it is the right first object: if
co-evolution does not appear here, it will not be identifiable in a
noisier range.
